\section{Cutting a matching: budgets and forced arches}
\label{sec:cut}

\subsection{The asymmetric cut bijection}

Fix one matching $P\in\NC_n$, a cut position $0\le L\le2n$ and $w=h_P(L)$; put
$R=2n-L$. The \emph{left inward word} of $P$ is the prefix of $w_P$ of length
$L$. The \emph{right inward word} is the suffix of length $R$ read from the right
end towards the cut with $U$ and $D$ exchanged. Both are \emph{ballot words}
(nonnegative prefix heights) ending at height $w$, and their numbers of $D$
letters, the \emph{down budgets}, are
\begin{equation}
\label{eq:budgets}
 A_L=\frac{L-w}{2},\qquad A_R=\frac{R-w}{2}.
\end{equation}
In particular $0\le w\le\min(L,R)$ and $w\equiv L\equiv R\pmod2$.

\begin{lemma}[Asymmetric cut bijection]
\label{lem:cut}
\lean{Meanders.FirstCrossing.asymmetricCutEquiv, Meanders.FirstCrossing.BallotHalf.down_budget, Meanders.FirstCrossing.through_pair_ordinal}
\leanok
\uses{lem:partition}
For every $L$ and $w$, the matchings $P\in\NC_n$ with $h_P(L)=w$ are in
bijection with pairs of ballot words of lengths $L$ and $R$, both ending at
height $w$. Under this bijection the internal arches of each half are the stack
matches of its word, and the $w$ \emph{through arches} join the $i$-th oldest
unmatched left opening to the $i$-th oldest unmatched right opening
(oldest meaning nearest the outer end, i.e.\ farthest from the cut) for
$i=1,\dots,w$.
\end{lemma}

\begin{proof}
\leanok
Cutting $P$ gives two such words. Conversely, given two words, match each $D$ to
the top of its stack within its own half; the unmatched left openings
$a_1<\dots<a_w$ and unmatched right openings $b_1>\dots>b_w$ (in physical
coordinates, so that $a_1$ and $b_1$ are the oldest) must be joined by through
arches. An internal arch cannot enclose an unmatched opening of its own half,
since closing it would have removed that opening from the stack first, so
internal arches never cross through arches. Through arches are pairwise nested
exactly when $a_i$ is joined to $b_i$: any other assignment has $i<j$ with
$a_i<a_j<b_{\sigma(i)}<b_{\sigma(j)}$, a crossing. Hence there is exactly one
noncrossing completion, and the two maps are mutually inverse.
\end{proof}

Both halves are read \emph{inward}, from the outer ends of the line towards the
cut. The reflection of the right half exchanges opening and closing roles but
does not change the matching or its owner. Nothing in the lemma relates $L$ to
$R$: the two halves may have very different lengths and budgets.

\subsection{The forced-prefix lemma}

Consider one inward ballot word of length $m$ with down budget $A$ and final
height $w=m-2A$. After a feasible prefix of length $i$ containing $c$ letters
$D$, the height is $h=i-2c$, the remaining length is $\ell=m-i$ and the
remaining down budget is $r=A-c$, with $0\le r\le\ell$. Number the currently
unmatched openings from oldest to newest.

\begin{lemma}[Exact forced prefix]
\label{lem:forced}
\lean{Meanders.FirstCrossing.forcedPrefix_exact, Meanders.FirstCrossing.forcedOpenings_card}
\leanok
\uses{lem:cut}
The set of currently unmatched openings that remain unmatched under
\emph{every} continuation of length $\ell$ with exactly $r$ letters $D$ and
nonnegative heights is precisely the oldest
\[
 f=\max(0,\,h-r)=\max(0,\,i-A-c)
\]
openings. Moreover $i-A-c$ increases by one at a $U$ letter and is unchanged
at a $D$ letter, so $f$ is nondecreasing along the word.
\end{lemma}

\begin{proof}
\leanok
A continuation contains $r$ letters $D$, so the height never drops below
$h-r$ and the oldest $\max(0,h-r)$ openings survive every continuation. For the
converse exhibit continuations that close every newer opening. If $h>r$ take
$D^{r}U^{\ell-r}$; if $h\le r$ take $D^{h}(UD)^{r-h}U^{\ell-2r+h}$. Both have
length $\ell$, use exactly $r$ letters $D$, stay nonnegative, end at height
$h+\ell-2r=w\ge0$, and close all but the oldest $f$ openings. Under a $U$ the
quantity $i-A-c$ grows by one; under a $D$ both $i$ and $c$ grow, so it is
unchanged, and $f=\max(0,i-A-c)$ cannot decrease.
\end{proof}

We call the oldest $f$ openings \emph{emitted}. Emitted openings can only be
closed by a through arch; the remaining $\rho=h-f=\min(h,A-c)$ newest openings
form a \emph{reservoir} that the word's own $D$ letters consume in stack order (\cref{fig:forced}).

\begin{figure}[t]
\centering
\begin{tikzpicture}[x=1cm,y=0.5cm,>=Stealth,every node/.style={font=\small}]
 \foreach \k/\lab in {0/1,1/2,2/3,3/4,4/5} {
   \draw (0,\k) rectangle (1.2,\k+0.9);
   \node at (0.6,\k+0.45) {$\lab$};
 }
 \node[anchor=east] at (-0.2,0.45) {oldest};
 \node[anchor=east] at (-0.2,4.45) {newest};
 \draw[decorate,decoration={brace,amplitude=4pt}] (1.4,4.9) -- (1.4,3.0)
   node[midway,anchor=west,xshift=6pt] {reservoir: newest $r=2$, may still close inside the half};
 \draw[decorate,decoration={brace,amplitude=4pt}] (1.4,2.9) -- (1.4,0)
   node[midway,anchor=west,xshift=6pt] {emitted: oldest $f=h-r=3$, forced through the cut};
 \node[anchor=north] at (0.6,-0.3) {$h=5$ open, $r=A-c=2$ downs left};
\end{tikzpicture}
\caption{\Cref{lem:forced} on one inward half: with $h=5$ current openings
and $r=2$ remaining down steps, at most the two newest openings can still be
closed inside the half, so the three oldest are emitted. Emitted openings
await their through partner in a FIFO; the reservoir is consumed in stack
order.}
\label{fig:forced}
\end{figure}

\begin{corollary}[Forced FIFO joins]
\label{cor:fifo}
\lean{Meanders.FirstCrossing.forcedFIFO_sound, Meanders.FirstCrossing.emitted_next_between}
\leanok
\uses{lem:cut, lem:forced}
Let both inward words of one matching be read in any interleaved order. As soon
as the $i$-th oldest opening has been emitted on both sides, the through arch
joining these two openings is present in every matching compatible with the
two prefixes and the budgets \eqref{eq:budgets}. Hence the join may be
performed immediately, and the two openings may be forgotten; the ordinal $i$ is
never affected by later letters.
\end{corollary}

\begin{proof}
\leanok
By \cref{lem:forced} an emitted opening is unmatched in every completion of its
half, and by \cref{lem:cut} the unmatched openings are joined by common oldest-first
ordinal. Newer openings cannot change the ordinal of an older one.
\end{proof}

Emission is cumulative: $f$ counts emitted openings including those already
joined and forgotten. At the end of a half, $c=A$, the reservoir is empty and
$f=w$, so all $w$ through arches have been joined once both halves are finished.
